\documentclass[11pt,a4paper]{article}

\usepackage[utf8]{inputenc}
\usepackage[T1]{fontenc}
\usepackage{lmodern}
\usepackage{amsmath,amssymb,amsfonts}
\usepackage{booktabs}
\usepackage{hyperref}
\usepackage[margin=1in]{geometry}
\usepackage{microtype}
\usepackage{xcolor}

\hypersetup{
  colorlinks=true,
  linkcolor=blue!60!black,
  citecolor=blue!60!black,
  urlcolor=blue!60!black,
}

\title{\textbf{Supplementary Material: The Standard Model from a Torus Knot}}
\author{James P.\ Galasyn and Claude Th\'eodore}
\date{}

\begin{document}
\maketitle

%----------------------------------------------------------------------
\section{S1. Self-Energy Derivation and \texorpdfstring{$\alpha$}{alpha} Emergence}
\label{sec:S1}

\subsection{S1.1 Torus knot parameterization}

A $(p,q)$ torus knot on a torus with major radius $R$ and minor radius $r$ is parameterized by $\lambda \in [0, 2\pi)$:
\begin{equation}
  \mathbf{r}(\lambda) = \begin{pmatrix} (R + r\cos q\lambda)\cos p\lambda \\ (R + r\cos q\lambda)\sin p\lambda \\ r\sin q\lambda \end{pmatrix}.
\end{equation}

The total path length $L = \int_0^{2\pi} |d\mathbf{r}/d\lambda|\,d\lambda$. For $(2,1)$ with $r \ll R$: $L \approx 4\pi R$.

\subsection{S1.2 Null condition and circulation energy}

The photon satisfies $ds^2 = 0$, giving $|\mathbf{v}| = c$ identically. The circulation energy is:
\begin{equation}
  E_{\text{circ}} = \frac{2\pi\hbar c}{L}.
\end{equation}

\subsection{S1.3 Electromagnetic self-energy}

Current $I = ec/L$. Magnetic self-energy from Neumann inductance~\cite{jackson1999}:
\begin{equation}
  U_{\text{mag}} = \frac{1}{2}\mu_0 R\left[\ln\frac{8R}{r} - 2\right]\left(\frac{ec}{L}\right)^{\!2}.
\end{equation}

For $v = c$: $U_{\text{elec}} = U_{\text{mag}}$ (transverse field equality). Total:
\begin{equation}
  \frac{U_{\text{EM}}}{E_{\text{circ}}} = \frac{\alpha[\ln(8R/r) - 2]}{\pi}.
\end{equation}

The fine-structure constant $\alpha = e^2/(4\pi\varepsilon_0\hbar c)$ appears automatically as the ratio of electromagnetic coupling to quantum of action.

\subsection{S1.4 Self-consistency: \texorpdfstring{$r/R = \alpha$}{r/R = alpha}}

The tube radius $r$ equals the transverse electromagnetic scale, giving $r/R = \alpha$. With this constraint, $E_{\text{total}}(R) = mc^2$ determines the torus size: $R_e = 194.2$~fm for the electron.

%----------------------------------------------------------------------
\section{S2. Self-Consistent Radii}
\label{sec:S2}

Bisection search on $R \in [10^{-20}, 10^{-8}]$~m, 60 iterations, ${\sim}15$-digit precision.

\begin{table}[h]
\centering
\begin{tabular}{@{}lccc@{}}
\toprule
Particle & Mass (MeV) & $R$ (fm) & $r$ (fm) \\
\midrule
$e$ & 0.5110 & 194.20 & 1.417 \\
$\mu$ & 105.66 & 0.939 & 0.00685 \\
$\tau$ & 1776.9 & 0.0558 & 0.000407 \\
\bottomrule
\end{tabular}
\caption{Self-consistent torus radii for the three charged leptons.}
\label{tab:radii}
\end{table}

Energy breakdown (electron): $E_{\text{circ}} = 0.5086$~MeV (99.5\%), $U_{\text{EM}} = 0.0024$~MeV (0.5\%).

%----------------------------------------------------------------------
\section{S3. Angular Momentum and Spin-\texorpdfstring{$\frac{1}{2}$}{1/2}}
\label{sec:S3}

The $(2,1)$ knot winds twice toroidally for every poloidal circuit. Under $2\pi$ coordinate rotation, the photon completes half its path --- the signature of spin-$\tfrac{1}{2}$. The time-averaged angular momentum $\langle L_z \rangle = (E/c) \int(xv_y - yv_x)\,ds / \int ds$ approaches $\hbar/2$ at the self-consistent aspect ratio.

%----------------------------------------------------------------------
\section{S4. Topology Uniqueness}
\label{sec:S4}

The torus (genus~1) is the simplest closed orientable surface supporting non-contractible loops. The $(2,1)$ knot is the simplest torus knot producing spin-$\tfrac{1}{2}$ with broken $p\leftrightarrow q$ symmetry.

\paragraph{Terminological note.} In standard knot theory, a $(p,q)$ torus knot with $p=1$ or $q=1$ is the unknot --- it is topologically trivial when embedded in $\mathbb{R}^3$. We use ``torus knot'' to refer to the closed geodesic on the torus surface, whose geometric properties (winding number, path length, self-energy) are the relevant features of the framework, independent of its embedding in $\mathbb{R}^3$. The $(2,1)$ curve is non-trivial as a geodesic on the torus even though it is trivial as a knot in ambient space.

The candidates are:
\begin{itemize}
  \item $(1,0)$: integer spin (boson)
  \item $(1,1)$: spin-$\tfrac{1}{2}$ but $p=q$ symmetric
  \item \textbf{$(2,1)$: spin-$\frac{1}{2}$, minimal asymmetric fermionic knot}
\end{itemize}

%----------------------------------------------------------------------
\section{S5. Koide Angle and Lepton Masses}
\label{sec:S5}

\subsection{S5.1 The Koide formula}

\begin{equation}
  Q = \frac{\sum m_i}{\left(\sum\sqrt{m_i}\right)^2} = \frac{2}{3}.
\end{equation}

Holds to 0.0009\% with PDG masses~\cite{pdg2022}. $Q = 2/3$ is automatic in the parameterization $\sqrt{m_i} = (S/3)(1 + \sqrt{2}\cos(\theta_K + 2\pi i/3))$; the content is the angle $\theta_K$.

\subsection{S5.2 Geometric derivation}

\begin{equation}
  \theta_K = \frac{p}{k}\left(\pi + \frac{q}{k}\right) = \frac{2}{3}\left(\pi + \frac{1}{3}\right) = \frac{6\pi + 2}{9}.
\end{equation}

Decomposition:
\begin{itemize}
  \item $2\pi/3$: $\mathbb{Z}_3$ generation symmetry from $k=3$ harmonic modes
  \item $2/9 = pq/k^2$: toroidal--poloidal coupling through $k\times k$ interaction matrix
\end{itemize}

Geometric value: $2.316617$~rad. Fitted value: $2.316616$~rad. Agreement: 0.00005\%.

\subsection{S5.3 Why \texorpdfstring{$\sqrt{m}$}{sqrt(m)} appears}

Mass $\propto 1/R$ in the torus model. Wavefunction normalization on the torus goes as $\sqrt{1/R} \propto \sqrt{m}$. Generation mixing involves overlap integrals of wavefunctions on tori of different sizes, naturally producing $\sqrt{m}$ ratios.

\subsection{S5.4 Lepton predictions}

With $m_e$ as input and $\theta_K = (6\pi+2)/9$:
\begin{itemize}
  \item $m_\mu = 105.659$~MeV (measured 105.658, 0.001\%)
  \item $m_\tau = 1777.0$~MeV (measured $1776.86 \pm 0.12$, within $1\sigma$)
\end{itemize}

%----------------------------------------------------------------------
\section{S6. Quark Koide Extension}
\label{sec:S6}

\subsection{S6.1 Generalized angle formula}

\begin{equation}
  \theta = \frac{6\pi + \dfrac{2}{1 + 3|q_{\text{em}}|}}{9}.
\end{equation}

\begin{table}[h]
\centering
\begin{tabular}{@{}ccccc@{}}
\toprule
Sector & $|q_{\text{em}}|$ & $\theta$ & $\Delta$ predicted & Error \\
\midrule
Leptons & 0 & $(6\pi+2)/9$ & 2.000 & 0.00\% \\
Down & $1/3$ & $(6\pi+1)/9$ & 1.000 & 0.85\% \\
Up & $2/3$ & $(6\pi+2/3)/9$ & 0.667 & 0.41\% \\
\bottomrule
\end{tabular}
\caption{Generalized Koide angle by charge sector.}
\label{tab:quark-koide}
\end{table}

Physical interpretation: the $k=3$ aspect ratio reduces the Koide angle. As $|q_{\text{em}}|$ increases (stronger coupling to the $k=3$ cavity modes), the angle approaches $2\pi/3$.

\subsection{S6.2 Hierarchy parameter}

\begin{equation}
  B^2 = 2(1 + C_F q_{\text{em}}^2), \quad C_F = \frac{k^2 - 1}{2k} = \frac{4}{3}.
\end{equation}

\begin{table}[h]
\centering
\begin{tabular}{@{}lccc@{}}
\toprule
Sector & $B^2$ predicted & $B^2$ fitted & Error \\
\midrule
Leptons & 2.000 & 2.000 & exact \\
Down & 2.296 & 2.389 & 3.9\% \\
Up & 3.185 & 3.094 & 2.9\% \\
\bottomrule
\end{tabular}
\caption{Hierarchy parameter by charge sector.}
\label{tab:hierarchy}
\end{table}

The SU($k$) Casimir $C_F$ appears because the $k=3$ harmonic mode interactions enhance mass splitting.

\subsection{S6.3 Scale invariance}

The Koide ratio $Q$ is invariant under $m_i \to \lambda m_i$, so QCD mass running (approximately flavor-universal at leading order) does not alter the angle or hierarchy measurements.

%----------------------------------------------------------------------
\section{S7. Anchor Mass Formulas}
\label{sec:S7}

\subsection{S7.1 Tube scale}

\begin{equation}
  \Lambda_{\text{tube}} = \frac{\hbar c}{\alpha R_p} = 255.67~\text{GeV}.
\end{equation}

\subsection{S7.2 Top quark --- geometric coupling \texorpdfstring{($\alpha^0$)}{(alpha0)}}

\begin{equation}
  m_t = \frac{p}{k}\Lambda = \frac{2}{3} \times 255{,}670 = 170{,}447~\text{MeV} \quad (\text{PDG: } 172{,}760,\ -1.3\%).
\end{equation}

The top quark couples directly to the tube energy. The coefficient $p/k = 2/3$ is the toroidal winding divided by the aspect ratio.

\subsection{S7.3 Bottom quark --- knot geodesic \texorpdfstring{$\times$}{x} EM \texorpdfstring{($\alpha^1$)}{(alpha1)}}

\begin{equation}
  m_b = \sqrt{p^2+q^2}\;\alpha\;\Lambda = \sqrt{5} \times \frac{\Lambda}{137.036} = 4{,}172~\text{MeV} \quad (\text{PDG: } 4{,}180,\ -0.2\%).
\end{equation}

The factor $\sqrt{5} = \sqrt{p^2+q^2}$ is the geodesic length of the $(2,1)$ torus knot on the unit covering space.

\subsection{S7.4 Tau lepton --- surface area / group order \texorpdfstring{$\times$}{x} EM \texorpdfstring{($\alpha^1$)}{(alpha1)}}

\begin{equation}
  m_\tau = \frac{p^2(p^2+q^2)}{k(2k+1)}\;\alpha\;\Lambda = \frac{20}{21}\alpha\Lambda = 1{,}776.9~\text{MeV} \quad (\text{PDG: } 1{,}776.86,\ +0.001\%).
\end{equation}

\begin{itemize}
  \item Numerator: $20 = p^2(p^2+q^2) = 4\times 5$ (toroidal area $\times$ knot metric)
  \item Denominator: $21 = k(2k+1) = 3\times 7$ (aspect ratio $\times$ extended group dimension)
\end{itemize}

This is one of the most precise predictions in the model (6~ppm). The factor $2k+1 = 7$ corresponds to the dimension of the fundamental representation of $\mathrm{SO}(2k+1) = \mathrm{SO}(7)$; its appearance in the lepton anchor mass warrants further investigation.

\subsection{S7.5 Inter-generation ratios}

\begin{align}
  m_c/m_t &\approx \alpha \quad (0.7\%\ \text{error}), \\
  m_s/m_b &\approx 3\alpha \quad (2.1\%\ \text{error}).
\end{align}

The charm-to-top ratio IS the fine-structure constant. The strange-to-bottom ratio is $\alpha \times k$.

%----------------------------------------------------------------------
\section{S8. Electroweak Sector}
\label{sec:S8}

\subsection{S8.1 Mode counting on the torus}

We conjecture that the $(2,1)$ torus knot admits 13 independent electromagnetic modes (a rigorous derivation from the spectral geometry is deferred to future work):
\begin{itemize}
  \item $p^2 = 4$ toroidal modes
  \item $q^2 = 1$ poloidal mode
  \item $p^2 q^2 = 4$ mixed (toroidal--poloidal) modes
  \item $p^2 + q^2 = 5$ knot-metric modes
  \item Minus 1 zero mode
  \item Total: $4 + 1 + 4 + 5 - 1 = 13$
\end{itemize}

Of these, 3 couple to the poloidal (weak isospin) sector, giving:
\begin{equation}
  \sin^2\theta_W = \frac{3}{13} \quad (\text{PDG: } 0.23122,\ 0.19\%\ \text{error}).
\end{equation}

\subsection{S8.2 Strong coupling from winding enhancement}

The strong coupling enhancement $p^4 = 16$ is observed to connect the electromagnetic and strong scales:
\begin{equation}
  \alpha_s = p^4 \alpha = 16\alpha = \frac{16}{137.036} = 0.11676 \quad (\text{PDG: } 0.1179 \pm 0.0009).
\end{equation}

A rigorous derivation of the $p^4$ enhancement from the torus winding geometry is an open problem. We note that $p^4 = (p^2)^2$ and $p^2 = 4$ counts the toroidal modes, suggesting the enhancement arises from a squared toroidal mode coupling.

\subsection{S8.3 Higgs boson mass}

The Higgs is the breathing mode (radial oscillation) of the torus:
\begin{equation}
  m_H = \left(\frac{q}{p} - \alpha\right)\Lambda_{\text{tube}}.
\end{equation}

Base frequency $q/p = 1/2$ (poloidal-to-toroidal ratio), with vacuum polarization correction $-\alpha$:
\begin{equation}
  m_H = (0.5 - 0.007297)\times 255.67 = 125.97~\text{GeV} \quad (\text{PDG: } 125.10 \pm 0.14).
\end{equation}

\subsection{S8.4 Vacuum expectation value}

\begin{equation}
  v = (1 - (p^2+q^2)\alpha)\Lambda_{\text{tube}} = (1 - 5\alpha)\Lambda_{\text{tube}}.
\end{equation}
\begin{equation}
  v = (1 - 0.03649)\times 255.67 = 246.34~\text{GeV} \quad (\text{PDG: } 246.22).
\end{equation}

The correction $5\alpha = (p^2+q^2)\alpha$ involves the knot metric. This is one of the most precise predictions (0.049\%).

\subsection{S8.5 W and Z boson masses}

From $v$ and $\sin^2\theta_W$ via standard electroweak relations:
\begin{itemize}
  \item $m_W = 80.4$~GeV (PDG 80.37, 0.0\%)
  \item $m_Z = 91.6$~GeV (PDG 91.19, 0.5\%)
\end{itemize}

%----------------------------------------------------------------------
\section{S9. CKM Matrix Derivation}
\label{sec:S9}

\subsection{S9.1 Cabibbo angle from Koide mismatch}

The up-type and down-type Koide angles differ by:
\begin{equation}
  \Delta\theta_{ud} = \theta_d - \theta_u = \frac{1}{27}~\text{rad}.
\end{equation}

This mismatch, amplified by the factor $2k = 6$, gives the Cabibbo angle:
\begin{equation}
  \theta_C = 6\Delta\theta_{ud} = \frac{6}{27} = \frac{2}{9}~\text{rad} = 12.73^\circ.
\end{equation}
\begin{equation}
  V_{us} = \sin\frac{2}{9} = 0.2214 \quad (\text{PDG: } 0.2243,\ 1.3\%).
\end{equation}

Compare with the Gatto relation: $\sin\theta_C \approx \sqrt{m_d/m_s} = 0.2237$.

\subsection{S9.2 \texorpdfstring{$V_{cb}$}{Vcb} from hierarchy mismatch}

\begin{equation}
  V_{cb} = \frac{|\Delta B|}{p^2+q^2} = \frac{|B_d^2 - B_u^2|/B_{\text{avg}}}{5} = 0.0429 \quad (\text{PDG: } 0.0408).
\end{equation}

The denominator is the knot metric $p^2+q^2 = 5$. The numerator measures the Koide hierarchy mismatch between charge sectors.

\subsection{S9.3 \texorpdfstring{$V_{ub}$}{Vub} from two-generation product}

\begin{equation}
  V_{ub} = \frac{p}{p^2+q^2}\times V_{us}\times V_{cb} = \frac{2}{5}\times 0.2214\times 0.0429 = 0.00376.
\end{equation}

PDG: 0.00382, error $-1.5\%$. The suppression factor $p/(p^2+q^2) = 2/5$ means two-generation mixing is the product of one-generation mixings, reduced by the knot ratio.

\subsection{S9.4 CP phase as deficit angle}

\begin{equation}
  \delta_{CP}^{\text{CKM}} = \pi - p = \pi - 2 = 1.14159~\text{rad} = 65.39^\circ.
\end{equation}

PDG: $1.144 \pm 0.027$~rad ($0.1\sigma$ agreement). CP violation is the geometric deficit: the $(2,1)$ winding $p = 2$ falls short of a half-turn $\pi$ by exactly $\pi - 2$.

%----------------------------------------------------------------------
\section{S10. PMNS Matrix Derivation}
\label{sec:S10}

\subsection{S10.1 Solar angle from mode counting}

\begin{equation}
  \sin^2\theta_{12} = \frac{p^2}{13} = \frac{4}{13} = 0.30769 \quad (\text{PDG: } 0.307).
\end{equation}

The same 13-mode denominator as $\sin^2\theta_W = 3/13$. The numerator $p^2 = 4$ counts toroidal modes coupling to the solar neutrino channel.

\paragraph{Consistency check:} $\sin^2\theta_W + \sin^2\theta_{12} = 3/13 + 4/13 = 7/13$. The shared denominator is a non-trivial structural prediction.

\subsection{S10.2 Atmospheric angle from knot metric}

\begin{equation}
  \sin^2\theta_{23} = \frac{p^2+q^2}{k^2} = \frac{5}{9} = 0.5556 \quad (\text{PDG: } 0.561).
\end{equation}

Numerator: knot metric $p^2+q^2 = 5$. Denominator: $k^2 = 9$.

\subsection{S10.3 Reactor angle from phase mismatch}

The neutrino Koide angle $\theta_\nu$ is determined by matching the mass-squared ratio $R = \Delta m^2_{31}/\Delta m^2_{21} = 33.83$. The reactor angle is:
\begin{equation}
  \sin^2\theta_{13} = \frac{(\theta_\nu - \theta_K)^2}{3} = 0.0218 \quad (\text{PDG: } 0.0220,\ 0.6\%).
\end{equation}

\subsection{S10.4 CP phase and mass ordering}

$\delta_{CP}^\nu = \pi$ ($180^\circ$), consistent with measured $177^\circ \pm 20^\circ$. Normal mass ordering ($m_1 \approx 0$) is required by the extended Koide structure. Both are testable at JUNO~\cite{juno2016} and Hyper-Kamiokande.

%----------------------------------------------------------------------
\section{S11. Neutrino Mass Scale}
\label{sec:S11}

\subsection{S11.1 Seesaw formula}

\begin{equation}
  A_\nu^2 = \frac{\alpha_W}{\pi}\frac{m_e^2}{\Lambda_{\text{tube}}},
\end{equation}
where $\alpha_W = \alpha/\sin^2\theta_W$. Using $\sin^2\theta_W = 3/13$: $A_\nu^2 = 0.01028$~eV (from data: 0.00989~eV, 3.9\% error).

\subsection{S11.2 Predictions}

\begin{table}[h]
\centering
\begin{tabular}{@{}lc@{}}
\toprule
State & Mass (meV) \\
\midrule
$m_1$ & 0.36 \\
$m_2$ & 8.66 \\
$m_3$ & 50.3 \\
$\Sigma m_\nu$ & 59.4 \\
\bottomrule
\end{tabular}
\caption{Predicted neutrino masses.}
\label{tab:neutrino-masses}
\end{table}

DESI+Planck bound: $\Sigma m_\nu < 64$~meV~\cite{desi2024}. Oscillation floor: 59.0~meV. Prediction: 59.4~meV --- sits precisely between these bounds.

%----------------------------------------------------------------------
\section{S12. Input Counting and Parameter Independence}
\label{sec:S12}

\subsection{S12.1 One primary input}

The model requires one primary measured input:
\begin{enumerate}
  \item $m_e = 0.51100$~MeV --- the electron mass (sets the absolute energy scale)
\end{enumerate}

The three quantities previously treated as additional inputs are now derivable from the topological integers $(p{=}2,\, q{=}1,\, k{=}R/r{=}3)$:
\begin{itemize}
  \item \textbf{$\alpha$}: Skilton's formula $\alpha^{-1} = \sqrt{137^2 + \pi^2} = 137.036016$ (0.12~ppm), where $137 = p_s^2 + q_s^2$ with generators $p_s = 2k + p^2 + q^2 = 11$ and $q_s = p^2 = 4$ --- these are NWT quantum numbers. See Paper~3~[ref].
  \item \textbf{$m_\mu$}: The Koide angle $\theta_K = (6\pi+2)/9$ is a function of $(p,q,k)$, so given $m_e$, the Koide relation determines $m_\mu = 105.658$~MeV (0.001\% error). See Paper~3~[ref].
  \item \textbf{$\Lambda_{\text{tube}}$}: The Lenz formula $m_p/m_e = 6\pi^5$ (0.002\% error) is structurally explained as $2k \times \pi^{p^2+q^2}$. Given $m_e$, this yields $m_p \to R_p \to \Lambda_{\text{tube}}$. See Paper~3~[ref].
\end{itemize}

Dimensional analysis requires at least one dimensionful input --- $m_e$ sets the absolute scale from which all other masses follow.

\subsection{S12.2 Three topological integers}

\begin{enumerate}
  \item $p = 2$ --- toroidal winding (spin-$\tfrac{1}{2}$ requirement)
  \item $q = 1$ --- poloidal winding (simplest closure)
  \item $k = R/r = 3$ --- torus aspect ratio (the lowest value admitting a Pythagorean resonance $(3,4,5)$, making the proton the unique stable baryon; quarks are harmonic modes of the $k{=}3$ cavity)
\end{enumerate}

\subsection{S12.3 Derived quantities}

From these four numbers (1 measured $+$ 3 topological), all 23 predictions follow through a derivation chain:
\begin{enumerate}
  \item $(p,q,k) \to \alpha$: via Skilton generators $p_s = 2k + p^2 + q^2 = 11$, $q_s = p^2 = 4$, giving $\alpha^{-1} = \sqrt{137^2 + \pi^2}$
  \item $(p,q,k) + m_e \to m_\mu$: via Koide angle $\theta_K = (6\pi+2)/9$ with topology-determined parameters
  \item $(p,q,k) + m_e \to m_p \to \Lambda_{\text{tube}}$: via Lenz formula $m_p/m_e = 6\pi^5 = 2k \times \pi^{p^2+q^2}$
  \item $\alpha + m_e + m_\mu + \Lambda_{\text{tube}} \to$ 23 SM parameters: via all the formulas of this paper
\end{enumerate}
\begin{itemize}
  \item $\theta_K = (6\pi+2)/9$ from $(p,q,k)$
  \item $\theta(q_{\text{em}}) = (6\pi + 2/(1+3|q|))/9$ from charge
  \item All coefficients: $2/3$, $\sqrt{5}$, $20/21$, $3/13$, $16$, $2/9$, $\pi-2$, etc.\ --- pure functions of $(p,q,k)$
\end{itemize}

See Paper~3~[ref] for the complete derivation of steps 1--3.

\subsection{S12.4 What is NOT an input}

\begin{itemize}
  \item $\sin^2\theta_W$, $\alpha_s$ --- predicted
  \item $v$, $m_H$, $m_W$, $m_Z$ --- predicted
  \item All CKM and PMNS parameters --- predicted
  \item 7 of 9 fermion masses --- predicted
  \item Neutrino masses and ordering --- predicted
\end{itemize}

%----------------------------------------------------------------------
\section{S13. Comparison with Other Approaches}
\label{sec:S13}

\begin{table}[h]
\centering
\begin{tabular}{@{}lccc@{}}
\toprule
Framework & SM params explained & Free params & Predictive? \\
\midrule
Standard Model & 0 (all input) & 26 & No \\
SU(5) GUT~\cite{georgi1974} & Gauge unification & ${\sim}24$ & Partially \\
SO(10) GUT & Gauge + partial Yukawa & ${\sim}15$--20 & Partially \\
$A_4$ flavor symmetry~\cite{ma2001} & Neutrino mixing & ${\sim}8$ & TBM only \\
String landscape~\cite{bousso2000} & None specific & $10^{500}$ vacua & No \\
\textbf{Torus knot (this work)} & \textbf{23} & \textbf{1 + 3 integers} & \textbf{Yes} \\
\bottomrule
\end{tabular}
\caption{Comparison of parameter-reduction frameworks.}
\label{tab:comparison}
\end{table}

%----------------------------------------------------------------------
\section{S14. Computational Verification}
\label{sec:S14}

All predictions are reproducible from the open-source simulation code:
\begin{verbatim}
python3 -m simulations.nwt --koide      # Lepton Koide analysis
python3 -m simulations.nwt --quarks     # Quark masses + EW/CKM/PMNS
python3 -m simulations.nwt --neutrino   # Neutrino masses and mixing
python3 -m simulations.nwt --self-energy # alpha emergence
python3 -m simulations.nwt --find-radii # Self-consistent torus radii
\end{verbatim}

No fitting routines or optimization are used. All predictions are computed from closed-form expressions evaluated at the geometric angle.

%----------------------------------------------------------------------
\begin{thebibliography}{10}

\bibitem{jackson1999}
J.~D.~Jackson,
\emph{Classical Electrodynamics}, 3rd ed.\ (Wiley, 1999).

\bibitem{pdg2022}
Particle Data Group, R.~L.~Workman \emph{et al.},
``Review of particle physics,''
\emph{Prog.\ Theor.\ Exp.\ Phys.}\ \textbf{2022}, 083C01 (2022); 2024 update.

\bibitem{juno2016}
F.~An \emph{et al.} (JUNO Collaboration),
``Neutrino physics with JUNO,''
\emph{J.\ Phys.\ G}\ \textbf{43}, 030401 (2016).

\bibitem{desi2024}
DESI Collaboration,
``DESI 2024 VI: Constraints on models of dark energy, neutrino mass, and modified gravity,''
arXiv:2404.03002 (2024).

\bibitem{georgi1974}
H.~Georgi and S.~L.~Glashow,
``Unity of all elementary-particle forces,''
\emph{Phys.\ Rev.\ Lett.}\ \textbf{32}, 438 (1974).

\bibitem{ma2001}
E.~Ma and G.~Rajasekaran,
``Softly broken $A_4$ symmetry for nearly degenerate neutrino masses,''
\emph{Phys.\ Rev.\ D}\ \textbf{64}, 113012 (2001).

\bibitem{bousso2000}
R.~Bousso and J.~Polchinski,
``Quantization of four-form fluxes and dynamical neutralization of the cosmological constant,''
\emph{J.\ High Energy Phys.}\ \textbf{2000}, 006 (2000).

\end{thebibliography}

\end{document}
